\section{Einleitung}
\label{sec:einleitung}

Deutschland ist nicht als isoliertes Stromsystem zu verstehen. Erzeugung,
Verbrauch und Handel sind in einen europäischen Verbund eingebettet, während
die tatsächlich gemessenen physikalischen Flüsse den elektrotechnischen
Eigenschaften des Netzes folgen. Handelsmenge und physischer Grenzfluss können
deshalb für dasselbe Länderpaar voneinander abweichen. Gleichzeitig verändern
sich erneuerbare Erzeugung und Day-Ahead-Preis stündlich. Wer Deutschlands Rolle
als Importeur oder Exporteur untersuchen möchte, muss daher räumliche,
zeitliche und mehrdimensionale Zusammenhänge gemeinsam betrachten.

Das Zielproblem dieses Projekts ist die explorative Analyse dieser
Zusammenhänge in einer einzigen Webanwendung. Eine blo{\ss}e Tabelle beantwortet
zwar einzelne Wertfragen, erschwert aber das Erkennen von Zeitmustern,
Richtungswechseln und länderübergreifenden Unterschieden. Umgekehrt würde eine
einzige aggregierte Kurve die Struktur der Partnerländer verdecken. Daraus
ergibt sich die Forschungsfrage:

\begin{quote}
\itshape
Wann ist Deutschland im betrachteten Zeitraum Stromimporteur oder
Stromexporteur, mit welchen Ländern findet der physische Austausch statt und
wie hängen diese Flüsse zeitlich mit Erzeugungsmix und Day-Ahead-Preis zusammen?
\end{quote}

Die Formulierung ist bewusst explorativ. Die Anwendung soll Muster sichtbar
machen und konkrete Stunden vergleichbar machen; sie soll keine Kausalmodelle,
Netzlastflussrechnung oder Herkunftsnachweise für einzelne Strommengen
ersetzen. Insbesondere bedeutet zeitliche Gleichläufigkeit nicht, dass eine
Erzeugungsart einen bestimmten Grenzfluss verursacht.

Das Vorgehen orientiert sich am verschachtelten Prozessmodell von Munzner:
Zuerst werden Anwendungssituation und Aufgaben beschrieben, danach folgen
Daten- und Aufgabenabstraktion, visuelle Kodierung und schlie{\ss}lich die
algorithmische Umsetzung \cite{munzner2009nested}. So wird vermieden, zuerst
leicht implementierbare Diagramme auszuwählen und die Fragestellung im
Nachhinein an diese anzupassen.

\subsection{Anwendungshintergrund}
\label{sec:hintergrund}

Im europäischen Verbundnetz verbinden Grenzkuppelleitungen benachbarte
Gebotszonen. Ein physischer Fluss beschreibt den gemessenen realen
Leistungsaustausch über eine solche Grenze. Ein Handelsfluss ist dagegen ein
kommerzielles Ergebnis gekoppelter Märkte. Physischer und kommerzieller Fluss
stimmen im Allgemeinen nicht überein, weil elektrische Leistung nicht dem
vertraglich gedachten kürzesten Weg folgt, sondern sich entsprechend der
Netzimpedanzen auf Leitungen verteilt \cite{smardFlows}. Die Anwendung stellt
beide Grö{\ss}en deshalb nicht gleich: Physische Flüsse bestimmen die grafische
Richtung und Stärke, Handelswerte werden als zusätzlicher Detailwert im
Tooltip gezeigt.

Die betrachtete deutsche Erzeugung ist ein Leistungsmix aus erneuerbaren und
konventionellen Technologien. Solar- und Windleistung variieren stark mit
Tageszeit und Wetter. Kohle- und Gaserzeugung können sich ebenfalls ändern,
unterliegen aber anderen technischen und wirtschaftlichen Randbedingungen.
Der Day-Ahead-Preis für die Gebotszone Deutschland--Luxemburg ist ein
Marktpreis pro Megawattstunde. Er kann negativ werden, wenn in einer Stunde ein
hohes Angebot auf eine vergleichsweise geringe Nachfrage und begrenzte
Flexibilität trifft. Die Visualisierung zeigt die zeitliche Gleichläufigkeit,
nicht die vollständigen Ursachen des Preises.

Energy-Charts stellt solche Daten durch das Fraunhofer-Institut für Solare
Energiesysteme ISE bereit. Für dieses Projekt liegen Kopien als PostgreSQL-
Views im Schema \code{energycharts} auf dem Seminarserver vor. Die offizielle
API-Dokumentation definiert die Einheiten: \code{total\_power} in MW,
grenzüberschreitende physische Flüsse (CBPF) und Handel (CBET) in GW sowie den
Preis in EUR/MWh \cite{energychartsOpenapi}. Positive CBPF-/CBET-Werte bedeuten
Import, negative Werte Export. Diese Definition ist die Grundlage sämtlicher
Legenden und Vorzeichenangaben.

\subsection{Zielgruppen}
\label{sec:zielgruppen}

Primäre Zielgruppe sind Studierende und fachlich interessierte Personen, die
Grundbegriffe wie Leistung, Erzeugungsart, Import und Export kennen, aber keine
spezielle Netzberechnungssoftware bedienen. Sie möchten Fragen wie
``Mit welchem Land floss in dieser Stunde besonders viel Leistung?'' oder
``Tritt ein negativer Preis gleichzeitig mit hoher erneuerbarer Erzeugung und
bestimmten Exportmustern auf?'' ohne SQL-Abfrage beantworten.

Eine zweite Zielgruppe sind Lehrende und energiewirtschaftlich Vorgebildete,
die den Datensatz zur Demonstration oder für eine erste Hypothesenbildung
nutzen. Von ihnen kann erwartet werden, dass sie Zeitreihen, positive und
negative Achsen sowie divergende Farbskalen lesen können. Für beide Gruppen
bleiben Einheiten, Vorzeichen und der Unterschied zwischen Physik und Handel
sichtbar. Fachbegriffe werden im Bericht erklärt; die Anwendung setzt kein
Wissen über Tabellennamen oder PostgREST voraus.

Die zentralen Informationsbedürfnisse sind:

\begin{itemize}
  \item Überblick über Richtung und Stärke aller Länderbeziehungen zu einer
        Stunde,
  \item Erkennen zeitlicher Veränderungen von Erzeugungsmix und ausgewähltem
        Länderfluss,
  \item schnelles Auffinden von länder- und zeitbezogenen Mustern in vielen
        Einzelwerten,
  \item Ablesen absoluter Werte und relativer Erzeugungsanteile bei Bedarf,
  \item Nachvollziehbarkeit der verwendeten Einheiten und Transformationen.
\end{itemize}

\subsection{Überblick und Beiträge}
\label{sec:beitraege}

Die Anwendung kombiniert einen gerichteten radialen Netzwerkgraphen, eine
gestapelte Zeitreihe mit separater Flussachse und eine Pixelmatrix. Alle drei
Ansichten verwenden denselben Partner- und Zeitindex. Ein Klick auf ein Land,
eine Stunde oder eine Matrixzelle aktualisiert die übrigen Ansichten. Der
Datensatz umfasst den gesamten Mai 2025 in stündlicher Auflösung.

Die Beiträge des Projekts sind:

\begin{enumerate}
  \item eine auf Anwendungsaufgaben zurückgeführte Auswahl dreier
        komplementärer Visualisierungstechniken aus den Bereichen Graphen,
        Zeitreihen und pixelorientierte Verfahren;
  \item eine reproduzierbare Datenpipeline von vier realen PostgreSQL-Views zu
        einem sicheren, statisch auslieferbaren HTTP-Datensatz;
  \item eine reine Elm-/SVG-Implementierung mit koordiniertem
        Auswahlzustand, mehreren Zeitfenstern und Details-on-Demand;
  \item eine fallbezogene Analyse des Mai 2025 einschlie{\ss}lich
        Datenqualitäts-, Build- und Interaktionstests.
\end{enumerate}

Die öffentliche Projekthistorie, Quelltexte, Daten und der Bericht befinden
sich im \projectlink{Hochschul-GitLab-Projekt}.
